% ======================================================================
\section{Complexity analysis of our own code}\label{sec:complexity}
% ======================================================================
\paragraph{One GA generation} A generation creates $\mu$ offspring. Per child: two tournaments $O(1)$,
OX $O(n)$ (marker array), mutation $O(n)$, Split $O(n)$ with the deque version ($O(nL)$ with Bellman;
$O(Kn)$ more in the rare fleet-limited case), the edge-set hash $O(n)$, clone look-up $O(1)$ expected, and
re-insertion into the sorted population $O(\mu)$. Hence
\begin{equation}
T_{\text{GA-gen}}=\Theta\big(\mu\,(n+\mu)\big)=\Theta(\mu n)\ \ (\mu\le n),\qquad
T_{\text{GA-gen}}^{\text{Bellman}}=\Theta(\mu nL),\qquad
T_{\text{MA-gen}}=T_{\text{GA-gen}}+p_{ls}\,\mu\cdot I_{LS}\cdot O(nk),
\end{equation}
where $I_{LS}$ is the number of local-search passes (data dependent; it grows with $n$ because a random child
needs more improving moves). Memory: $\Theta(\mu n)$.

\paragraph{One TS iteration} The neighbourhood has at most $7nk+n$ moves, each evaluated in $O(1)$ by
\cref{eq:delta,eq:delta2} plus an $O(1)$ tabu test; applying the chosen move rebuilds at most two routes,
$O(L)$; finding an empty slot is $O(K)$ and the penalty update $O(1)$. Copying a new incumbent costs $O(n)$
but happens rarely. So
\begin{equation}
T_{\text{TS-it}}=\Theta(nk+L+K)=\Theta(nk),\qquad\text{and }\Theta(n^2)\text{ for the full neighbourhood }k=n-1 .
\end{equation}

\paragraph{What delta evaluation saves} Re-evaluating a neighbour from scratch costs $\Theta(n)$ (all
routes) or $\Theta(L)$ if only the two modified routes are re-summed; one iteration would then cost
$\Theta(n^2k)$ or $\Theta(nkL)$ instead of $\Theta(nk)$ --- a factor $\Theta(n)$, respectively $\Theta(L)$.
Measured per move (\texttt{bench-eval}, averaged over 5 passes over the whole neighbourhood), delta
evaluation costs 1.0--1.3\,ns up to $n=800$ (2--4\,ns beyond, when the matrix leaves the cache), route-level
re-evaluation 35--70\,ns and full re-evaluation 49\,ns at $n=25$ up to 2.4\,$\mu$s at $n=3200$: delta
evaluation is 41$\times$ to 548$\times$ faster and the gap widens with $n$, as predicted. All three methods
return identical deltas.

\begin{figure}[t]
\centering
\includegraphics[width=\linewidth]{\figdir/complexity.pdf}
\caption{Empirical validation. (a) Median time per TS iteration and per GA/MA generation on synthetic
instances with constant route length (and the real X instances), log--log; the number at each line end is the
fitted slope. (b) For nine quantities, the exponent predicted by the analysis (hollow) versus the fitted
slope with its 95\,\% confidence interval (filled). (c) Split for $n=1000$ and growing route length $L$:
Bellman's $\Theta(nL)$ against the deque's $\Theta(n)$.}
\label{fig:complexity}
\end{figure}

\paragraph{Empirical validation} We timed both algorithms on synthetic instances with
$n=25\ldots3200$ and constant expected route length ($Q=155$, about 10 customers per route), each value the
median of three runs of a fixed number of iterations on an otherwise idle machine; initialisation cancels out
because we time two run lengths and take the difference. \Cref{fig:complexity} and \cref{tab:slopes} compare
the slopes of the log--log regressions with the predictions.
\begin{table}[h]
\centering\small
\caption{Predicted exponents versus least-squares slopes of $\log T$ on $\log n$ (or $\log L$).}
\label{tab:slopes}
\begin{tabular}{lcc}
\toprule
Measured quantity & predicted exponent & fitted slope (95\% CI) \\
\midrule
TS iteration, granular $k=20$ & 1 & $1.08\pm0.06$ \\
TS iteration, full neighbourhood & 2 & $2.12\pm0.07$ \\
TS iteration on X instances & 1 & $0.97\pm0.04$ \\
GA generation, linear Split & 1 & $0.84\pm0.10$ \\
MA generation (incl.\ local search) & -- & $1.38\pm0.08$ \\
one move, $O(1)$ delta & 0 & $0.01\pm0.09$ \\
one move, route re-evaluation & 0 ($L$ fixed) & $0.07\pm0.05$ \\
one move, full re-evaluation & 1 & $0.78\pm0.20$ \\
Bellman Split vs.\ $L$ ($n$ fixed) & 1 & $0.57\pm0.12$ \\
linear Split vs.\ $L$ ($n$ fixed) & 0 & $-0.09\pm0.04$ \\
\bottomrule
\end{tabular}

\end{table}

\paragraph{Explaining the differences} (i) \emph{TS} matches its prediction ($1.08$, $2.12$; $0.97$ on the
real X instances) up to $n=1600$. At $n=3200$ the time per iteration jumps by $5.9\times$ instead of $2\times$:
the $(n{+}1)^2$ distance and tabu matrices (2$\times$82\,MB) no longer fit in the processor caches, so each
matrix look-up becomes a memory access --- a constant of the RAM model that silently changes. (ii) The \emph{GA}
slope $0.84$ is below 1 because a per-child constant (allocations, random numbers) dominates for small $n$;
between $n=800$ and $3200$ the local slope is $0.99$. (iii) \emph{Full re-evaluation} ($0.78$) and
\emph{Bellman versus $L$} ($0.57$) are sums $a+bn$ resp.\ $a'n+b'nL$ whose constant part still matters at the
measured sizes; their local slopes approach 1 at the large end, while the deque Split stays flat ($-0.09$: fewer
routes have to be built when $L$ grows) and becomes $37\times$ faster than Bellman at $L=500$. (iv) The MA has no
clean exponent because the number of local-search moves per child is itself instance dependent --- a reminder that
``cost per iteration'' and ``iterations needed'' are different questions.
The Clarke--Wright baseline itself takes 25--220\,$\mu$s on the six main instances.
